\chapter{Compresibilidad y ondas}\label{ch:08}
\objetivos{Derivar la velocidad del sonido y las relaciones isentrópicas de un gas perfecto; relacionar Mach con propagación de información y cambio de área; distinguir ondas débiles y choques; calcular saltos normales y la geometría ideal de un cono de Mach.}

Una corriente gaseosa no necesita ser muy rápida en metros por segundo para mostrar compresibilidad: importa la velocidad relativa a la propagación de perturbaciones y la magnitud de las variaciones de presión. El número de Mach organiza esta comparación. Los materiales docentes presentan velocidad del sonido, cono de Mach y conductos compresibles \cite{TAE3}; aquí desarrollamos los balances que permiten interpretar sus fórmulas y corregir signos mediante argumentos físicos.

\section{Compresibilidad y perturbación acústica}
La compresibilidad mide cómo responde la densidad a cambios de presión bajo un proceso especificado. Una compresión lenta con intercambio térmico no produce la misma relación que una perturbación rápida adiabática. Por eso la derivada que define velocidad del sonido debe indicar la restricción termodinámica:
\begin{equation}\label{eq:08-sonido}
 a^2=\left(\frac{\partial p}{\partial\rho}\right)_s.
\end{equation}
El subíndice $s$ expresa entropía constante para una perturbación débil ideal, reversible y adiabática.

Podemos obtener esta relación linealizando las ecuaciones unidimensionales alrededor de un gas uniforme en reposo: $\rho=\rho_0+\rho'$, $p=p_0+p'$ y $u=u'$, con perturbaciones pequeñas. Al despreciar productos de perturbaciones,
\begin{equation}
 \frac{\partial\rho'}{\partial t}+\rho_0\frac{\partial u'}{\partial x}=0,
 \qquad \rho_0\frac{\partial u'}{\partial t}=-\frac{\partial p'}{\partial x}.
\end{equation}
Con $p'=a^2\rho'$, derivamos una ecuación y sustituimos en la otra:
\begin{equation}
 \frac{\partial^2p'}{\partial t^2}-a^2\frac{\partial^2p'}{\partial x^2}=0.
\end{equation}
Las soluciones se propagan a velocidad $a$ respecto al gas. Si este se mueve uniformemente, la propagación axial respecto al observador es $U+a$ o $U-a$. Esta última velocidad muestra por qué una corriente supersónica no transmite pequeñas perturbaciones aguas arriba en el modelo unidimensional.

Para gas ideal caloríficamente perfecto y proceso isentrópico, $p/\rho^\gamma$ es constante, de donde
\begin{equation}
 a=\sqrt{\gamma p/\rho}=\sqrt{\gamma RT},\qquad \gamma=c_p/c_v.
\end{equation}
La presión no aparece de forma independiente al expresar el resultado con $T$: a temperatura y composición iguales, el modelo da igual velocidad del sonido aunque cambie la densidad. En gases reales, a temperaturas altas o cerca de cambios importantes de composición, esa simplificación necesita revisión.

\section{Mach y selección de régimen}
El número de Mach es $M=U/a$ y utiliza la velocidad relativa al aire local. Si $M<1$ el flujo es subsónico; si $M>1$, supersónico. El régimen transónico incluye coexistencia de regiones subsónicas y supersónicas, de modo que un Mach de corriente libre inferior a uno puede acompañarse de choques locales sobre un perfil. El término hipersónico suele emplearse a partir de valores de Mach del orden de cinco, pero su interés físico incluye calentamiento e interacciones que no comienzan todas en un límite único.

Para variaciones suaves de velocidad a bajo Mach, la variación relativa de densidad suele escalar como $M^2$. La regla de usar un modelo incompresible por debajo de aproximadamente $M=0.3$ es una orientación de error, no una ley que prohíba cambios de densidad por calentamiento a velocidades pequeñas. Un flujo lento con gran aporte de calor puede tener variaciones considerables de densidad sin deberse a la aceleración.

La información acústica y la materia se mueven de forma distinta. Una onda no transporta necesariamente una masa de aire a la misma velocidad a la que avanza su frente. En el límite lineal, las partículas oscilan alrededor de su estado medio mientras la perturbación se propaga. Este hecho impide interpretar el sonido como una corriente de aire emitida continuamente por el vehículo.

\section{Temperatura y presión de remanso}
En un flujo estacionario adiabático sin trabajo de eje ni cambios apreciables de potencial, la energía da $h+U^2/2=h_0$. Para gas caloríficamente perfecto,
\begin{equation}\label{eq:08-T0}
 T_0=T+\frac{U^2}{2c_p},\qquad
 \frac{T_0}{T}=1+\frac{\gamma-1}{2}M^2.
\end{equation}
La temperatura de remanso $T_0$ representa el estado alcanzado al frenar idealmente la corriente sin pérdida de energía total. Un sensor real puede recuperar solo una fracción del incremento de temperatura debido a la geometría y transferencia de calor.

Si el frenado es además isentrópico,
\begin{equation}\label{eq:08-p0}
 \frac{p_0}{p}=\left(1+\frac{\gamma-1}{2}M^2\right)^{\gamma/(\gamma-1)},\qquad
 \frac{\rho_0}{\rho}=\left(1+\frac{\gamma-1}{2}M^2\right)^{1/(\gamma-1)}.
\end{equation}
La energía por sí sola no implica conservación de $p_0$: una irreversibilidad puede disminuir presión de remanso manteniendo temperatura de remanso. Las relaciones isentrópicas no atraviesan una onda de choque \cite{NASA-BGA}.

Estas expresiones permiten invertir una medida de impacto subsónica ideal:
\begin{equation}
 M=\left\{\frac{2}{\gamma-1}\left[\left(\frac{p_0}{p}\right)^{(\gamma-1)/\gamma}-1\right]\right\}^{1/2}.
\end{equation}
A bajo Mach, una expansión de Taylor recupera $p_0-p\simeq\rho U^2/2$. Así se ve que Bernoulli incompresible es un límite, no una fórmula independiente que deba aplicarse a cualquier velocidad.

\section{Conductos: derivación de área--velocidad}
Consideremos flujo cuasiunidimensional, estacionario, inviscido e isentrópico en un conducto. La continuidad diferencial proporciona
\begin{equation}
 \frac{d\rho}{\rho}+\frac{dU}{U}+\frac{dA}{A}=0.
\end{equation}
Euler sobre la corriente da $dp=-\rho U\,dU$ y la relación acústica isentrópica $dp=a^2d\rho$. Por tanto, $d\rho/\rho=-M^2dU/U$. Al sustituir,
\begin{equation}\label{eq:08-area}
 \frac{dA}{A}=(M^2-1)\frac{dU}{U}.
\end{equation}
En subsónico, un convergente acelera y reduce presión; un divergente frena y aumenta presión. En supersónico ocurre lo opuesto: un divergente acelera y reduce presión, mientras un convergente frena y aumenta presión. El signo de presión se comprueba siempre con $dp=-\rho U\,dU$. Acelerar y aumentar presión a la vez no es compatible con estas hipótesis.

Para pasar suavemente de subsónico a supersónico hace falta una sección extrema con $dA=0$ cuando $M=1$. En la tobera convergente--divergente ideal es la garganta mínima. La geometría es una condición necesaria para esa solución, pero no suficiente: las presiones de contorno deben permitirla. Una tobera con parte divergente no genera automáticamente flujo supersónico en toda condición de alimentación.

Integrando junto con las relaciones isentrópicas se obtiene
\begin{equation}
 \frac{A}{A^*}=\frac{1}{M}
 \left[\frac{2}{\gamma+1}\left(1+\frac{\gamma-1}{2}M^2\right)\right]^{(\gamma+1)/[2(\gamma-1)]},
\end{equation}
con $A^*$ área sónica correspondiente al mismo caudal y estado de remanso. Para $A/A^*>1$ existen una rama subsónica y otra supersónica. El contexto y las condiciones de contorno seleccionan la rama; resolver una ecuación algebraica sin examinarlas deja indeterminado el problema físico.

\section{Estrangulamiento y presión de salida}
Con $p_0$, $T_0$ y área de garganta fijos, el caudal ideal crece al reducir la presión de descarga hasta alcanzar Mach uno en la sección mínima. Después se dice que el flujo está estrangulado: una reducción adicional de presión aguas abajo no aumenta el caudal dentro de ese modelo. Para una garganta sónica,
\begin{equation}
 \dot m=\frac{p_0A^*}{\sqrt{T_0}}\sqrt{\frac{\gamma}{R}}
 \left(\frac{2}{\gamma+1}\right)^{(\gamma+1)/[2(\gamma-1)]}.
\end{equation}
La presión crítica cumple $p^*/p_0=[2/(\gamma+1)]^{\gamma/(\gamma-1)}$, aproximadamente $0.5283$ para $\gamma=1.4$. Esta relación corresponde al estado sónico, no a cualquier salida supersónica.

En una tobera real puede aparecer un choque, separación o expansión exterior si la presión de descarga no coincide con la de la solución ideal. El caudal estrangulado no significa que todo el interior sea supersónico ni que el empuje sea independiente del entorno. Las condiciones exteriores siguen afectando a la adaptación de presión y a la estructura de ondas.

\section{Choques: conservación con irreversibilidad}
Una onda de choque es una región delgada con cambios rápidos de presión, densidad, temperatura y velocidad. En un modelo macroscópico se trata como una discontinuidad y se aplican balances entre estados uniformes a ambos lados. Un choque normal es perpendicular a la corriente incidente. Para estado $1$ anterior y $2$ posterior,
\begin{align}
 \rho_1U_1&=\rho_2U_2,\\
 p_1+\rho_1U_1^2&=p_2+\rho_2U_2^2,\\
 c_pT_1+\tfrac12U_1^2&=c_pT_2+\tfrac12U_2^2.
\end{align}
La región se considera adiabática sin trabajo de eje; eso no la hace reversible. La entropía aumenta para el choque físico compresivo. Eliminando variables con gas ideal y $\gamma$ constante,
\begin{align}
 \frac{p_2}{p_1}&=1+\frac{2\gamma}{\gamma+1}(M_1^2-1),\label{eq:08-shock-p}\\
 \frac{\rho_2}{\rho_1}&=\frac{(\gamma+1)M_1^2}{2+(\gamma-1)M_1^2},\\
 M_2^2&=\frac{1+(\gamma-1)M_1^2/2}{\gamma M_1^2-(\gamma-1)/2},\qquad
 \frac{T_2}{T_1}=\frac{p_2/p_1}{\rho_2/\rho_1}.
\end{align}
Con $M_1>1$ el choque normal produce $M_2<1$, aumento de presión estática y temperatura y reducción de velocidad. $T_0$ se conserva y $p_0$ disminuye. En el límite $M_1\to1$ los saltos tienden a uno y se recupera una perturbación débil. A Mach muy alto, suponer calores específicos constantes puede dar temperaturas insuficientemente realistas.

Un choque oblicuo forma un ángulo $\beta$ con la corriente y se estudia mediante la componente normal $M_{n1}=M_1\sin\beta$. Su ángulo no es, en general, el ángulo de Mach. Una expansión supersónica suave se realiza mediante un abanico de ondas débiles; un choque de expansión violaría la condición de entropía en este contexto. Estas diferencias explican por qué no se deben dibujar todas las ondas como líneas equivalentes.

\section{Problema resuelto: estado tras un choque normal}
\begin{enumerate}
\item \textbf{Enunciado.} Una corriente a Mach dos atraviesa un choque normal estacionario. Calcular su estado posterior y comprobar conservación de energía y producción de entropía.
\item \textbf{Datos.} $M_1=2$, $p_1=24000\,\mathrm{Pa}$, $T_1=220\,\mathrm{K}$, $R=287.05287\,\mathrm{J\,kg^{-1}\,K^{-1}}$ y $\gamma=1.4$.
\item \textbf{Incógnitas.} $M_2$, $p_2$, $T_2$, velocidades y variación de entropía específica.
\item \textbf{Esquema verbal.} Flujo uniforme de izquierda a derecha que cruza una superficie normal sin cambio de área.
\item \textbf{Hipótesis.} Gas ideal caloríficamente perfecto, choque adiabático, sin trabajo de eje ni cambio de potencial; estados de equilibrio a ambos lados.
\item \textbf{Modelo.} Saltos normales obtenidos de masa, momento y energía, con criterio de entropía.
\item \textbf{Ecuaciones.} Relaciones anteriores y $\Delta s=c_p\ln(T_2/T_1)-R\ln(p_2/p_1)$, con $c_p=\gamma R/(\gamma-1)$.
\item \textbf{Despeje.} Se calculan primero razones de presión y densidad, después temperatura y Mach, y finalmente $U_i=M_i\sqrt{\gamma RT_i}$.
\item \textbf{Sustitución.} $p_2/p_1=4.5$, $\rho_2/\rho_1=8/3$ y $T_2/T_1=1.6875$.
\item \textbf{Resultado.} $p_2=108.0\,\mathrm{kPa}$, $T_2=371.25\,\mathrm{K}$, $M_2=0.57735$, $U_1=594.68\,\mathrm{m\,s^{-1}}$, $U_2=223.01\,\mathrm{m\,s^{-1}}$ y $\Delta s=93.95\,\mathrm{J\,kg^{-1}\,K^{-1}}$.
\item \textbf{Interpretación.} La corriente se frena y calienta, aumentando presión estática. La presión de remanso se reduce: $p_{02}/p_{01}=\exp(-\Delta s/R)=0.72087$.
\item \textbf{Comprobación independiente.} Continuidad da $U_2/U_1=\rho_1/\rho_2=3/8$, compatible con las velocidades. Además, $T_{01}=220(1+0.2\cdot4)=396\,\mathrm{K}$ y $T_{02}=371.25(1+0.2/3)=396\,\mathrm{K}$.
\item \textbf{Limitaciones.} No describe un choque oblicuo ni resuelve su espesor interno. A mayor temperatura deben revisarse calores específicos, vibración y composición.
\end{enumerate}

\section{Cono de Mach y retraso de un observador}
Una fuente puntual supersónica en movimiento rectilíneo uniforme emite perturbaciones débiles que forman una envolvente cónica. En un intervalo $\Delta t$, la fuente recorre $U\Delta t$ y una perturbación se extiende una distancia $a\Delta t$. La geometría de tangencia da
\begin{equation}
 \sin\mu=\frac{a\Delta t}{U\Delta t}=\frac1M,
\end{equation}
con semiángulo $\mu$ del cono. El cono representa el límite ideal de perturbaciones débiles; la firma de presión de una aeronave finita y los choques reales requieren información adicional.

\figura{mach}{Envolvente original de perturbaciones acústicas para una fuente a Mach 2. El semiángulo es de onda infinitesimal y no representa un choque finito.}

Si la fuente vuela a altura constante $h$ y el cronómetro comienza cuando pasa exactamente sobre un observador fijo, después de $t$ se ha alejado $Ut$. La llegada de la envolvente satisface $\tan\mu=h/(Ut)$. Combinando con $\sin\mu=1/M$,
\begin{equation}\label{eq:08-retraso}
 t=\frac{h}{a}\sqrt{1-\frac1{M^2}},\qquad
 M=\frac1{\sqrt{1-(at/h)^2}}.
\end{equation}
Para Mach supersónico finito, $0<t<h/a$. Un retraso inferior a $h/a$ es posible: la señal que llega fue emitida antes del paso por la vertical. Confundirla con una señal emitida en ese mismo instante conduciría a la relación incorrecta $t=h/a$.

Por ejemplo, con $h=1200\,\mathrm{m}$, $a=335\,\mathrm{m\,s^{-1}}$ y $t=2.5\,\mathrm{s}$, resulta $at/h=0.69792$, $M=1.3963$, $U\simeq467.8\,\mathrm{m\,s^{-1}}$ y $\mu\simeq45.74^\circ$. Como comprobación, $\tan\mu=h/(Ut)\simeq1.026$. Si se mantuviesen $h$ y $a$ pero se midiese $t=4\,\mathrm{s}$, no habría solución bajo estas hipótesis, porque $h/a=3.58\,\mathrm{s}$. No sería correcto aceptar una raíz imaginaria como Mach físico.

El modelo ignora viento, refracción por perfiles térmicos, maniobras y relieve. Además, al acercarse $at/h$ a uno, la inversión es muy sensible: pequeñas incertidumbres producen grandes cambios en Mach. Esa sensibilidad limita la utilidad de estimaciones basadas solo en un retraso cercano a su cota superior.

\section{Una comprobación transversal de los modelos}
Conviene ordenar las aproximaciones por sus hipótesis, sin suponer que una ecuación más corta sea siempre menos correcta. Bernoulli incompresible es adecuado para un frenado lento con pequeñas variaciones de densidad. La relación isentrópica de remanso permite densidad variable y excluye irreversibilidad. Las relaciones de choque admiten una discontinuidad irreversible y conservan los balances entre estados. Cada modelo responde a una pregunta distinta; elegirlo requiere identificar el proceso físico antes de sustituir números.

\section{Ejercicios propuestos}
\begin{enumerate}
\item \textbf{Básico.} Calcular $a$ en aire a $250\,\mathrm{K}$ y el Mach de una corriente a $200\,\mathrm{m\,s^{-1}}$.
\item \textbf{Intermedio.} Para $M=0.8$ y $T=260\,\mathrm{K}$, calcular $T_0$ y $p_0/p$ con $\gamma=1.4$.
\item \textbf{Examen.} Una fuente a $h=1500\,\mathrm{m}$ vuela a Mach $1.8$ con $a=320\,\mathrm{m\,s^{-1}}$. Hallar retraso y semiángulo de Mach, y comprobar que el retraso satisface la condición de existencia.
\item \textbf{Conceptual.} Explicar por qué un divergente puede acelerar un flujo o frenarlo y qué información necesita para escoger la opción.
\end{enumerate}
\section{Respuestas y orientación}
1. $a\simeq316.97\,\mathrm{m\,s^{-1}}$, $M\simeq0.631$. 2. $T_0=293.28\,\mathrm{K}$ y $p_0/p\simeq1.5243$. 3. $t\simeq3.898\,\mathrm{s}$, $\mu\simeq33.75^\circ$ y $h/a=4.6875\,\mathrm{s}$. 4. El signo de $M^2-1$ cambia la relación área--velocidad; se necesita conocer régimen y condiciones de contorno, y revisar pérdidas o choques.

\resumen{La velocidad del sonido procede de una perturbación isentrópica pequeña y Mach compara propagación de información y advección. Las relaciones de remanso y área requieren hipótesis específicas. Un choque conserva masa, momento y energía, pero aumenta entropía y pierde presión total. El ángulo de Mach describe ondas débiles y su geometría permite un retraso inferior al tiempo de propagación vertical desde el instante de sobrevuelo.}
\section{Autoevaluación}
\begin{enumerate}
\item ¿Por qué adiabático no implica isentrópico a través de un choque?
\item ¿Qué condición adicional a una garganta necesita una tobera para alcanzar la rama supersónica?
\item ¿Por qué el retraso del cono de Mach puede ser menor que $h/a$ sin violar propagación acústica?
\end{enumerate}
